\documentclass[12pt]{article}

\usepackage[a4paper, margin=1in]{geometry}
\usepackage{setspace}
\usepackage{titlesec}
\usepackage{fancyhdr}
\usepackage{xcolor}
\usepackage{ragged2e}

\setstretch{1.3}
\setlength{\parindent}{0pt}
\setlength{\parskip}{0.6em}
\AtBeginDocument{\RaggedRight}

\definecolor{HeadingBlue}{HTML}{1F4E79}
\titleformat{\section}{\Large\bfseries\color{HeadingBlue}}{}{0pt}{}
\titleformat{\subsection}{\large\bfseries\color{HeadingBlue}}{}{0pt}{}
\titlespacing*{\section}{0pt}{1.5em}{0.8em}
\titlespacing*{\subsection}{0pt}{1em}{0.6em}

\title{Islamic Studies Semester Project\\
\large The Core Divide: Understanding the Clash Between Secular Ideas and Islamic Beliefs}

\author{
Muhammad Hassaan Sajid (25L-2555)\\
Muhammad Ibrahim (25L-0716)\\
Ali Junaid (25L-0712)\\
Mehroz Mehdi (25L-0753)\\
Ali Hassan (25L-0768)\\
Aziz (25L-0808)
}

\date{}

\begin{document}

\maketitle

\newpage

% ================= MEMBER 1 =================
\newpage
\pagestyle{fancy}
\fancyhf{}
\fancyhead[R]{Muhammad Ibrahim (25L-0716)}
\fancyfoot[C]{\thepage}

\section*{The Clash of Truth: Where Do We Get Our Morals?}

\subsection*{Introduction}

Think about living in a society where laws regulating traffic change each and every day 
according to what becomes popular in social networks. That must be total chaos! However, 
when talking about an aspect which is far more serious than road safety regulations, that is our 
basic understanding of ethics, this is precisely how the modern world functions. In today’s 
world, ethics are considered to be as flexible as any other social phenomenon, and can be 
changed whenever new societal fashions arise. It leads us straight into a very important and 
unavoidable dilemma. 

"Do humans create their own rules, or do they come from God?" 

Modern secular culture tends to act like this question is already settled. Secularism puts a l its 
trust in human reason. Islam's response is not that reason is useless, just that it bends too 
easily, toward whatever is popular, toward our instincts, toward whoever holds power at the 
time. 

\subsection*{God's Word vs. Human Logic}

The secular claim is that reason alone is enough to build a fair society. The problem is that 
secular morality rests on changing opinions, not on anything fixed. Push on it and it moves. 

That became clear in a debate between Muslim speaker Hamza Tzortzis and physicist 
Lawrence Krauss. Tzortzis posed the "incest chalenge": is incest objectively wrong? Krauss, 
reasoning from a strictly secular and scientific view, admitted he could not cal it "absolutely" 
wrong, not if both people agreed and there was no genetic risk. The answer was honest. And 
that honesty was the problem. Once morality has nothing holding it in place other than human 
preference, even the most basic moral instincts are open to debate. 

Moroccan philosopher Taha Abderrahmane takes this further. Building ethics purely on human 
reason, with no room for Divine Revelation, does not free the soul. It empties it. Concepts like 
selflessness or loving your neighbor stop being moral truths and start being survival habits that 
happened to stick around. That is not nothing. But it is not a foundation. 

\subsection*{The Myth of the Neutral State}

The secular state presents itself as a neutral referee, managing different beliefs fairly without 
pushing any of its own. That is a convenient story, but not an accurate one. 

Legal scholar Wael Halaq argues the modern state functions more like a replacement god. It 
takes fu l control over laws and the use of force, insists religion stay private, and puts itself 
above every other loyalty. It does not share space with other ways of living. It pushes them out. 

Abdelwahab El-Messiri caled this the trap of "partial secularism." Separating religion from 
politics sounds reasonable at first. But the logic keeps going. It spreads into economics, 
culture, and private life, what he ca led "comprehensive secularism." At that point the secular 
state has not simply removed religion from public life. It has built its own version of one: 
founding documents treated as sacred texts, historical figures treated like prophets, national 
ceremonies replacing religious ones. It is not a neutral space. It is a different belief system 
wearing a neutral face. 

\subsection*{Changing Laws vs. Fixed Rules}

Islamic law holds certain Thawabit, fixed rules that do not shift based on political pressure or 
changing trends. The ban on Riba (exploitative interest) does not get rewritten because 
modern finance finds it inconvenient. 

Secular law has no equivalent. Written by people, based on agreement, it changes when the 
agreement changes. 

Under old English common law, a wife had almost no independent legal identity. That system 
was torn down and rebuilt around personal choice, with marriage becoming a flexible contract 
based on what two people currently want. 

The U.S. Supreme Court evaluates fundamental rights against "evolving standards of decency." 
In 1989, it ruled that executing 16-year-olds and menta ly disabled people was legal. About a 
decade later, as public opinion shifted, it reversed that ruling entirely. Same court, same 
document, opposite decision. The rule did not change. The feeling did. 

\subsection*{Conclusion}

If a moral system rewrites itself whenever public opinion shifts or a lobby pushes hard enough, 
it is not rea ly making moral claims. It is describing what people currently prefer and ca ling it 
justice. Whether Islam's fixed rules are the right answer is a real debate worth having. But the 
secular alternative, a morality that changes because we decide we want it to, is not obviously 
more reliable. It might just be more comfortable. 

% ================= MEMBER 2 =================
\newpage
\pagestyle{fancy}
\fancyhf{}
\fancyhead[R]{Mehroz Mehdi (25L-0753)}
\fancyfoot[C]{\thepage}

\section*{Arrival of secularism in Europe and spread to Muslim countries}

\subsection*{Introduction:}

Secularism in Europe developed gradually over centuries as a response to conflicts between religious authority and political power. In the medieval period, the Catholic Church strongly controlled society and governance, which created tension and criticism. The Renaissance encouraged human reason, while the Protestant Reformation weakened the unity of the Church. Later, wars like the Thirty Years’ War showed the dangers of mixing religion with politics. The Enlightenment further promoted freedom, tolerance, and rational thinking. As a result of these changes, the idea of separating religion from the state emerged, leading to the rise of secularism in Europe.

\subsection*{Main historical events:}

\textbf{1. The Renaissance (14th–17th century)}

\textbullet\ \textbf{Revival of Classical Knowledge}

The Renaissance marked a major revival of ancient Greek and Roman knowledge in fields such as philosophy, science, and politics. Scholars began studying thinkers like Aristotle and Plato, who emphasized logic, observation, and rational thinking. As a result, people slowly shifted away from relying solely on religious explanations and started exploring the world through reason and evidence. Another key feature of the Renaissance was the rise of humanism, which focused on human dignity, creativity, and intellectual potential. Education expanded beyond religious teachings to include literature, art, and science. This encouraged individuals to think independently and believe in the power of human reasoning.

\textbullet\ \textbf{Main Cause of Secularism (The rise of human reason over blind religious authority)}

This was the first major turning point where people began mentally separating knowledge and learning from religion, laying the foundation of secular thinking.

\textbf{2. The Protestant Reformation (16th century)}

\textbullet\ \textbf{Challenge to Church Authority}

The Protestant Reformation began when Martin Luther criticized the corruption of the Church, especially the practice of selling indulgences. He argued that salvation depends on faith rather than Church control. This directly challenged the authority and dominance of the Church. As a result of the Reformation, Europe was divided into Catholic and Protestant regions. The unity of Christianity broke down, and no single Church remained supreme. This weakened centralized religious power significantly. The Reformation also encouraged individuals to read religious texts themselves and form personal beliefs. Religion became more individual rather than controlled by a single institution.

\textbullet\ \textbf{Main Cause of Secularism (Breaking the monopoly of the Church over religion and authority)}

Once the Church lost its unified control, it could no longer dominate politics completely, opening the way for secular ideas.

\textbf{4. The Peace of Westphalia}

\textbullet\ \textbf{End of Religious Wars}

The Peace of Westphalia (1648) ended the Thirty Years’ War and reduced large-scale religious conflicts in Europe. It allowed rulers to choose their state’s religion without Church interference, strengthening political independence. As a result, the Church’s political influence declined, and states became more autonomous and self-governing.

\textbullet\ \textbf{Main Cause of Secularism (Recognition of state sovereignty over religious authority)}

This formally reduced the Church’s political power and established the idea that the state should operate independently.

\textbf{5. The Enlightenment (17th–18th century)}

\textbullet\ \textbf{Emphasis on Reason and Science}

The Enlightenment was an intellectual movement that emphasized reason, logic, and scientific thinking. Thinkers like John Locke and Voltaire promoted questioning authority and relying on rational thought instead of blind belief. They supported individual rights such as freedom of religion, speech, and thought, arguing that people should not be forced into a single religious system. Enlightenment philosophers also criticized the Church’s influence in politics and advocated a clear separation between religion and state.

\textbullet\ \textbf{Main Cause of Secularism (Philosophical support for separation of Church and state)}

This gave secularism a strong intellectual foundation and justified it as a necessary system.

\textbf{6. The French Revolution}

\textbullet\ \textbf{Attack on Church Privileges}

During the French Revolution, the Church’s privileges, wealth, and political influence were removed, and it lost its superior status. The state took control of key areas like education and administration, reducing the role of religion in governance. Religion was separated from state affairs, and people were recognized as equal citizens regardless of their faith. Laws were based on equality rather than religious authority.

\textbullet\ \textbf{Main Cause of Secularism (Practical removal of Church power from the state)}

This was the stage where secularism was fully implemented, turning ideas into reality.

\subsection*{Events that transferred it to Muslim states:}

\textbf{1. European Colonization (18th–20th century)}

\textbullet\ \textbf{Political Control by European Powers}

European empires like Britain and France took control of many Muslim regions (India, Egypt, Algeria, etc.). They introduced their own legal and administrative systems. Colonial rulers often replaced Islamic law (Sharia) in governance with European legal codes. Law and governance became separated from religion.

\textbullet\ \textbf{Main Impact}

Direct imposition of secular political and legal systems

This was the strongest external force spreading secularism.

\textbf{2. Decline of the Ottoman Empire}

\textbullet\ \textbf{Weakening of Islamic Political Unity}

The Ottoman Empire was one of the last major Islamic political powers. Its decline reduced unified Islamic governance. In the 19th century, the Ottomans introduced reforms inspired by Europe—modern laws, education, and administration. These reforms reduced the role of religion in state affairs.

\textbullet\ \textbf{Main Impact}

Adoption of European-style reforms within a Muslim empire

This made secular ideas enter from within, not just from outside.

\textbf{4. Western Education and Intellectual Influence}

\textbullet\ \textbf{Spread of European Ideas}

Muslim students studied in Europe and brought back ideas of democracy, nationalism, and secularism. Some Muslim intellectuals argued for modernization by adopting certain Western political ideas.

\textbullet\ \textbf{Main Impact}

Intellectual acceptance of secular political concepts

This made secularism seem “modern” and progressive.

\textbf{5. Formation of Modern Nation-States (20th century)}

End of Empires

After World War I, many Muslim regions became independent states. New states often adopted constitutions, legal systems, and governance structures based on European models.

\subsection*{Islamic historical model}

In classical Islamic history, the relationship between religion and governance was generally cooperative rather than confrontational.

\textbf{Role of Scholars (Ulama)}

Islamic scholars provided guidance on law, ethics, and religious interpretation, while rulers managed administration and governance. Both roles were distinct but connected.

\textbf{Integrated System}

In many Islamic empires (like the early Caliphates and later Muslim states), religious law and political authority worked together rather than competing. There was no long-term institutional “war” between religious authority and the state in the same way as Europe. This created a system where religion and governance were integrated rather than separated.

\subsection*{Main difference:}

\begin{itemize}
\item Europe developed secularism after a historical conflict between Church and State, where both competed for supremacy.
\item Islamic political tradition generally involved cooperation between scholars and rulers, not a structural conflict requiring separation.
\end{itemize}
\newpage
\pagestyle{fancy}
\fancyhf{}
\rhead{Muhammad Hassaan Sajid(25L-2555)}
\cfoot{\thepage}

\section*{Law and Power: Who Has the Right to Make the Rules?}

\textbf{Introduction}

Imagine a country where people can only practice their religion at home or in mosques. The government does not make laws based on religion and schools do not teach it. At first this might seem like a balance. For Islam, which says it is a complete way to live, this creates a big issue. The question is, can Islam really just be a part of life or does it have to be part of everything?

If Islam is a way of life, then how can it not be involved in making laws, teaching in schools, or running the country? This is a question that needs to be thought about:

\textit{``Can Islam survive if it is reduced to a private belief instead of a public system?''}

The big problem between secularism and Islam starts here. This issue is not about laws. It is about who gets to decide what is allowed and what is not. The question is whether Islam gets to define what is permissible and forbidden or something else does. The final authority in what is permitted and what is forbidden is what is really at stake.

\textbf{The Law of Heaven vs The Law of Men}

At the heart of secular democracy is a simple idea:

\begin{quote}
``Humans are the ultimate lawmakers.''
\end{quote}

Parliaments, constitutions, and public opinion together establish laws. The legal system allows for changes when society changes its opinions. This system relies on human reasoning and collective decision-making to define justice.

Islam takes a fundamentally different position. It accepts human reasoning to some extent but designates Divine Revelation (Wahi) as the ultimate source of authority. The Qur’an states:

\begin{quote}
``The rule is for none but Allah.'' (Surah Yusuf 12:40)
\end{quote}

This implies humans can implement laws but cannot overrule what has been revealed by Allah.

Another verse strengthens this principle:

\begin{quote}
``And whoever does not judge by what Allah has revealed; then it is those who are the disbelievers.'' (Surah Al-Ma’idah 5:44)
\end{quote}

The issue here is the choice between stability and flexibility. Human-made laws can change based on political systems, public opinion, and societal norms. What is acceptable today may not be acceptable tomorrow.

Islamic law, however, is based on permanent principles (Thawabit). These constants create a stable framework that remains unchanged despite social evolution.

The Qur’an emphasizes this completeness:

\begin{quote}
``This day I have perfected for you your religion and completed My favor upon you and have approved for you Islam as your religion.'' (Surah Al-Ma’idah 5:3)
\end{quote}

The tension is clear:

\textit{Should law follow human will, or should humans submit to a higher law?}

\textbf{Separating Professional and Personal}

Secularism promotes the idea that religion should remain personal and separate from law, politics, and governance. Individuals are free to believe what they want, but religion should not influence public systems.

However, Islam does not fit into this framework. Islam is not just a belief system; it is a complete way of life (Deen), covering worship, ethics, social relations, and governance.

The Qur’an commands:

\begin{quote}
``O you who believe! Enter into Islam completely (in total submission)...'' (Surah Al-Baqarah 2:208)
\end{quote}

Islam loses its completeness when reduced to personal practice only. The Qur’an also warns:

\begin{quote}
``Then do you believe in part of the Scripture and reject the rest?'' (Surah Al-Baqarah 2:85)
\end{quote}

This creates a contradiction:

\textit{A Muslim is asked to follow Islam completely, but the system they live in only allows partial practice.}

\textbf{Divided Justice}

In many modern Muslim societies, both Islamic and secular systems exist side by side. Islamic law is often limited to personal matters like marriage, divorce, and inheritance, while criminal law and financial systems are governed by secular frameworks.

For example:
\begin{itemize}
\item Family laws may follow Shariah principles
\item Financial systems may operate on interest (Riba)
\item Criminal justice may follow secular codes
\end{itemize}

This raises a key question:

\textit{If Divine law is truly just, why is it limited only to certain areas of life?}

The Qur’an states:

\begin{quote}
``And do not follow their desires away from what has come to you of the truth.'' (Surah Al-Ma’idah 5:48)
\end{quote}

And also:

\begin{quote}
``But no, by your Lord, they will not truly believe until they make you (O Prophet) judge in all disputes between them...'' (Surah An-Nisa 4:65)
\end{quote}

The debate over secularism and Islam ultimately revolves around authority: who defines truth and justice?

Secularism prioritizes human reasoning, while Islam prioritizes Divine guidance, using human intellect within its framework.

When both systems are combined without clarity, contradictions arise in law, governance, and ethics.

\textbf{Conclusion}

The real question is not whether religion should exist in society.

It is who has the final word in defining truth, justice, and law.
\end{document}
